% Part: lambda-calculus
% Chapter: syntax
% Section: abbreviated-syntax

\documentclass[../../../include/open-logic-section]{subfiles}

\begin{document}

\olfileid[fr]{lam}{syn}{abb}
\olsection{Syntaxe abrégée}

Les termes définis dans la \olref[trm]{defn:term}
sont parfois fastidieux à écrire. Il est donc
utile d'introduire une notation plus concise.
Il faut toutefois veiller à ce que les termes
ainsi notés restent de lecture unique. Une
convention courante, celle des \emph{termes
abrégés}, est la suivante.

\begin{enumerate}
\item Quand les parenthèses sont omises,
  l'application s'associe de gauche à droite.
  Par exemple, si $M$, $N$, $P$ et~$Q$ sont
  des termes, $MNPQ$ abrège $(((MN)P)Q)$.
\item De même, quand les parenthèses sont omises,
  l'abstraction lambda a la portée la plus
  large possible. Par exemple,
  $\lambd[x][MNP]$ se lit
  $(\lambd[x][((MN)P)])$.
\item Un lambda peut abstraire plusieurs
  variables. Par exemple,
  $\lambd[xyz][M]$ abrège
  $\lambd[x][\lambd[y][\lambd[z][M]]]$.
\end{enumerate}

Par exemple,
\[
\lambd[xy][xxyx \lambd[z][xz]]
\]
abrège
\[
(\lambd[x][(\lambd[y][((((xx)y)x)(\lambd[z][(xz)]))])]).
\]

\begin{prob}
Développer le terme abrégé
$\lambd[g][(\lambd[x][g (x x)])
  \lambd[x][g (x x)]]$.
\end{prob}

\end{document}
